\subsection{Definition}\label{subsec:generalized-reeb-orbits-polytopes-definition}

Let
\[
  K=\{x\in\mathbb{R}^4:\langle a_i,x\rangle\leq 1,\; i=1,\ldots,F\}
\]
be a full-dimensional polytope with \(0\in\operatorname{int}(K)\). We assume
that this \(H\)-representation is irredundant, so the facet corresponding to
\(a_i\) is
\[
  F_i=\{x\in K:\langle a_i,x\rangle=1\}.
\]
In this normalization the contact-normalized Reeb direction on \(F_i\) is
\begin{equation}\label{eq:generalized-reeb-facet-vector}
  R_i=2J_0a_i.
\end{equation}
Indeed \(R_i\) is tangent to \(F_i\), because
\(\langle a_i,J_0a_i\rangle=0\), and for \(x\in F_i\),
\[
  \lambda_0|_x(R_i)
  =\frac{1}{2}\langle -J_0R_i,x\rangle
  =\langle a_i,x\rangle
  =1.
\]

A \emph{generalized Reeb orbit} on \(K\) is a curve
\(\gamma\in W^{1,2}([0,T],\mathbb{R}^4)\) with \(T>0\),
\(\gamma(0)=\gamma(T)\), image contained in \(\partial K\), and
\begin{equation}\label{eq:generalized-reeb-inclusion}
  \dot{\gamma}(t)
  \in \operatorname{conv}\{R_i:\gamma(t)\in F_i\}
  \quad\text{for a.e. }t\in[0,T].
\end{equation}
This is the polytope form of the Hamiltonian inclusion
\(-J_0\dot{\gamma}\in\partial g_K^2(\gamma)\), translated into the
contact-normalized Reeb directions. Since every \(R_i\) in the active convex
hull satisfies \(\lambda_0(R_i)=1\) at \(\gamma(t)\), a generalized Reeb orbit
has
\[
  A(\gamma)=T.
\]

A generalized Reeb orbit is \emph{simple} if its periodic parametrization can
be started so that \(\dot{\gamma}\), after changing values at finitely many
breakpoints if necessary, is constant on finitely many intervals and each
constant value is one of the pure facet directions \(R_i\). In addition, each
facet direction occurs in one connected time interval, possibly not at all.
Thus a simple orbit uses pure facet directions rather than nontrivial convex
combinations, and it does not return to the same facet direction in separated
time intervals.

